
\subsection{Integrador duplo}

\subsubsection{init - configuracoes da simulacao}

dynamics of the problem



Omega: $\Big[1, 2\Big]$

Time step of 0.01 second.

Xref: $\Big[2, -1\Big]$

Tmin: 0.25

\begin{equation}
    Q = \left[ \begin{matrix} 1 & 0 \\ 0 & 1 \end{matrix} \right]
\end{equation}

Tpmax: 1

smax: 2

r: 1

\begin{equation}
    Ts = \left[ \begin{matrix} 0 & 2 & 3 & 5 & 6 \end{matrix} \right]
\end{equation}


$X_0$: $\Big[-0.5, -1\Big]$



The double integrator is a simple example where the matrices A and b define the state-space representation of a system.

\begin{equation}
    \begin{matrix}
        A_1 = \left[ \begin{matrix} 0 & 1 \\ 0 & 0 \\ \end{matrix} \right] &
        A_2 = \left[ \begin{matrix} 0 & 1 \\ 0 & 0 \\ \end{matrix} \right] \\
        \\
        b_1 = \left[ \begin{matrix} 0 \\ 1 \end{matrix} \right] &
        b_2 = \left[ \begin{matrix} 0 \\ 0 \end{matrix} \right]
    \end{matrix}
\end{equation}

The minimum switching time is set to 0.25.

The reference state vector is [2, -1], and the states can be interpreted as position and velocity of the system.

The variables Tpmax and smax are set to 1 and 2 respectively, but their specific roles are not clear from this context.

The variable r is set to 1, but its role is also not clear from this context.

The config structure seems to be a configuration for different operating modes of the system.
It includes fields for operation mode (Omega), associated control (ur), state-space matrices (A, b, C, D),
time step (tstep), reference state (xref), minimum time step (tmin), state weighting matrix (Q), maximum time period (Tpmax),
maximum value for s (smax), a variable r, time sequence (Ts), and initial state (x0).




\subsubsection{calculo da trajetoria (optimizacao)}

\subsubsection{construindo MPC (matrizes PHI e GAMMA)}

\subsubsection{rodando simulacao}

\subsubsection{resultados}